\chapter{Introduction}
\label{ch:intro}

\section{What this project is}
A \emph{compiler} is a program that translates text written by a human (source code) into something a machine can execute. Compilers are among the most intricate programs in computing because they must be \emph{exactly} right: a single wrong rewrite silently changes the meaning of every program that passes through it.

This project, called \nova{}, is a complete compiler-engineering platform. It is not one compiler but \emph{two}, plus the tools needed to build, test, debug and visualise them:

\begin{enumerate}
  \item \textbf{\nova{} the language and its compiler} --- a small, statically typed language (integers, floats, booleans, strings, arrays, functions) compiled either to a bytecode virtual machine or to real x86-64 machine code.
  \item \textbf{\novasql{}} --- a compiler for the SQL query language, together with an in-memory database engine that executes the compiled queries.
  \item \textbf{An assistant and a language server} --- tools that turn compiler diagnostics into verified fixes and bring the compiler into code editors.
  \item \textbf{A web studio} --- a React application whose 20 pages show every phase of both compilers using the \emph{real} compiler output, never mock data.
\end{enumerate}

\begin{plain}
Think of the compiler as a factory line. Raw material (source text) enters at one end; at each station it is reshaped --- split into words (lexing), arranged into sentences (parsing), checked for meaning (semantic analysis), simplified (optimization) --- and the finished product (machine code) leaves at the other end. This report walks along the whole line, explains what every station does, and shows how the stations were tested.
\end{plain}

\section{Motivation and goals}
Textbook compiler courses present each phase in isolation, usually on toy inputs. Real compilers are different in three ways that this project deliberately tackles:

\begin{itemize}
  \item \textbf{Phases must work together.} A type error found late, an optimization that breaks a loop, or a code generator that mishandles a corner case only shows up when all phases run on the same program.
  \item \textbf{Correctness is hard to see.} Compiler bugs are silent. The only convincing evidence of correctness is \emph{independent} evidence: a second implementation that computes the same answer by a different route.
  \item \textbf{Tooling matters.} Good diagnostics, a debugger, an editor integration and a way to \emph{look inside} the compiler are what make a compiler usable.
\end{itemize}

The goals were therefore: (1) implement every classical phase from scratch with no parser-generator or code-generator libraries; (2) support \emph{two} different kinds of languages (an imperative one and a declarative one) to show which techniques are universal; (3) verify everything by differential testing and fuzzing; (4) measure performance honestly; (5) make every internal step visible and understandable.

\section{Contributions}
\begin{enumerate}[label=\textbf{C\arabic*.}]
  \item A \textbf{canonical LR(1) parser generator} built from first principles (items, closure, goto, ACTION/GOTO tables) for a grammar with \LRProductions{} productions, producing \LRStates{} states without a single conflict, plus yacc-style panic-mode error recovery.
  \item A \textbf{typed front end} with scoped symbol tables, static type checking and edit-distance ``did you mean'' suggestions; \NumDiagCodes{} diagnostic codes with exact source ranges.
  \item A \textbf{middle end} with a three-address IR, CFG construction, dominators, liveness, reaching definitions, \emph{pruned SSA}, and seven optimization passes including \textbf{interprocedural inlining} with a call-graph recursion check.
  \item \textbf{Two back ends}: a stack-based \textbf{bytecode VM} with a time-travel debugger and static verifier, and a \textbf{native x86-64 code generator} with a \textbf{linear-scan register allocator}, a heap, strings and runtime traps.
  \item \textbf{\novasql{}}: parser, binder with \textbf{three-valued logic}, logical plan, rule-based optimizer (\textbf{predicate pushdown, hash joins, top-N}), expression bytecode and executor.
  \item A \textbf{compiler-verified assistant}: quick fixes and an auto-fixer that only offer a change after re-compiling it, rule-based runtime root-cause analysis, optional LLM integration whose output is re-checked, and an LSP server with code actions.
  \item A \textbf{verification framework}: a reference interpreter, a random program generator, a failure minimizer, native-versus-VM differential tests, a SQL-versus-SQLite fuzzer, mutation fuzzing and sanitizer builds --- \TotalTests{} automated tests in total.
  \item A \textbf{web platform}: a sandboxed REST service (\NumEndpoints{} endpoints) and a React/TypeScript studio.
\end{enumerate}

\section{The project at a glance}
\begin{table}[H]
\centering\small
\begin{tabular}{@{}lr@{\hspace{2.2cm}}lr@{}}\toprule
\textbf{Measure} & \textbf{Value} & \textbf{Measure} & \textbf{Value}\\ \midrule
Lines of code (all languages) & \TotalLOC{} & LR(1) states / conflicts & \LRStates{} / \LRConflicts{}\\
Compiler core + SQL (C++) & \CompilerLOC{} & Grammar productions & \LRProductions{}\\
\novasql{} engine (C++) & \SQLLOC{} & VM opcodes & \NumOpcodes{}\\
Automated tests & \TotalTests{} & Optimization passes & 7\\
\quad C++ / backend / frontend & \CppTests{} / \BackendTests{} / \FrontendTests{} & REST endpoints & \NumEndpoints{}\\
Example programs shipped & \NumExamples{} & Diagnostic codes & \NumDiagCodes{}\\
Git commits & \NumCommits{} & Studio pages & 19\\ \bottomrule
\end{tabular}
\caption{Key figures, generated from the repository (see Section~\ref{sec:reproduce}).}
\end{table}

\section{Repository tour}
Table~\ref{tab:loc} shows where the code lives. The compiler core is C++17 for speed and control; the web service is Python (FastAPI) because it mostly orchestrates processes and JSON; the studio is TypeScript/React.

\begin{table}[H]\centering\small
\begin{tabular}{@{}c@{}}\begin{tabular}{@{}lr@{}}\toprule
\textbf{Component} & \textbf{Lines}\\ \midrule
Lexer & 217\\
LR(1) parser generator and parser & 547\\
AST construction & 225\\
Semantic analysis & 482\\
IR generation & 333\\
Optimizer (incl. inlining) & 495\\
CFG / SSA / data-flow analyses & 650\\
Bytecode generator{,} verifier{,} x86-64 backend & 648\\
Bytecode VM and reference interpreter & 456\\
Pipeline{,} diagnostics{,} CLI{,} fuzzers & 740\\
NOVA-SQL compiler and engine & 1{,}605\\
C++ headers & 1{,}045\\
Native runtime (C) & 66\\
C++ unit tests & 1{,}612\\
Backend service (Python) & 1{,}811\\
Backend tests (Python) & 690\\
Frontend studio (TypeScript/React) & 3{,}763\\
Scripts (shell/Python) & 570\\
\midrule
\textbf{Total} & \textbf{15{,}955}\\ \bottomrule
\end{tabular}
\end{tabular}
\caption{Lines of code per component (measured by the report generator).}\label{tab:loc}
\end{table}

\begin{lstlisting}[style=tiny,caption={Top-level layout of the repository.},label={lst:layout}]
nova/
  compiler/     C++17: lexer, parser (LR(1) generator), ast, semantic, ir, optimizer, analysis,
                codegen (VM bytecode + x86-64), vm, sql/, tests/, runtime/ (C runtime for native code)
  backend/      FastAPI service (app/) and its pytest suite (tests/)
  frontend/     React + TypeScript + Vite + Tailwind + Monaco + React Flow (20 pages) + Vitest
  examples/     NOVA programs (valid, error, runtime-error) and SQL scripts
  benchmarks/   fib, primes and inlining workloads, with CPython equivalents
  editors/      VS Code client for the language server
  scripts/      dev, build, serve, test, sanitize, bench, native build, SQL fuzzer, LSP launcher
  docs/         specifications and design notes
  report/       this report (LaTeX) and the generator that produces its data
\end{lstlisting}

\section{How to read this report}
Every chapter starts with the idea in plain language, then gives the technical detail, then shows real output. Coloured boxes mark the different kinds of content:

\begin{plain}[In plain words]
An intuitive explanation for readers new to compilers.
\end{plain}
\begin{keybox}[Key idea]
The central technical insight of a section.
\end{keybox}
\begin{bugbox}[A real bug the tests caught]
A defect found by the project's own verification, how it was found and how it was fixed.
\end{bugbox}
\begin{warnbox}[Honest limitation]
Something the implementation deliberately does \emph{not} do, or has not been verified to do.
\end{warnbox}

\begin{table}[H]\centering\small
\begin{tabular}{@{}lp{0.62\textwidth}@{}}\toprule
\textbf{If you want to\ldots} & \textbf{read}\\ \midrule
understand the whole system quickly & Chapters~\ref{ch:intro}, \ref{ch:arch} and~\ref{ch:verify}\\
learn how a language is compiled & Chapters~\ref{ch:language}--\ref{ch:backends}\\
learn how SQL is compiled & Chapter~\ref{ch:sql}\\
see how correctness was established & Chapter~\ref{ch:verify} and the bug boxes throughout\\
run the project yourself & Appendix~\ref{app:manual}\\
look up a term & the Glossary at the end\\ \bottomrule
\end{tabular}
\caption{Reader paths.}
\end{table}

\section{Reproducing the numbers in this report}\label{sec:reproduce}
No number in this report was typed by hand. The script \code{report/tools/gen\_data.py} runs the real compiler, the test suites, the fuzzers and the benchmarks, and writes LaTeX fragments (tables, listings and macros such as \verb|\TotalTests|) that this document includes. Re-running the script and re-building the report therefore re-measures everything on the reader's machine.
